% Hafta 14 — Kontrol akışı: düzleştirmeden önce ve sonra
% Derleme: py -3.12 tools/latex/derle.py docs/week-14/assets/h14-10-cfg-once-sonra.tex
\documentclass[tikz,border=8pt]{standalone}
\usepackage{cen429-cizim}
\begin{document}
\begin{tikzpicture}[node distance=8mm and 16mm]
  \node[baslik] (b) at (0,0) {Kontrol akışı: düzleştirmeden önce ve sonra};
  \node[altbaslik, text width=150mm] at (b.south west) {aynı sonucu üreten iki program, çok farklı iki grafik};

  \node[iyikutu=62mm, below=16mm of b.south west, anchor=north west] (sol) {\kb{ÖNCE --- okunur}\\[2mm]
    [hesapla] $\rightarrow$ [karşılaştır]\\$\rightarrow$ [geç] / [red]\\birkaç blok\\akış gözle izlenir};
  \node[kotukutu=62mm, right=of sol] (sag) {\kb{SONRA --- düzleştirilmiş}\\[2mm]
    [switch] $\leftrightarrow$ \{çok sayıda case\}\\+ sahte durumlar\\çok blok\\sıra durum değişkeninde saklı};
  \draw[draw=cizgi, line width=1.6pt] ($(sol.north east)!0.5!(sag.north west)$) -- ($(sol.south east)!0.5!(sag.south west)$);

  \node[dip, text width=160mm, below=8mm of sol.south -| sag, anchor=north] {%
    Amaç: blok sayısını ve dallanmayı artırıp statik okumayı pahalı kılmak.};
\end{tikzpicture}
\end{document}
